% ===========================================================
% Глава 1. Анализ предметной области и существующих решений
% Подключить: \input{chapter1.tex}
% ===========================================================

\section{Анализ предметной области и существующих решений}

\subsection{Задача транскрибации многоязычной речи}

Задача, решаемая в рамках настоящей работы, может быть сформулирована следующим образом. На вход системе подаётся аудиофайл записи встречи, содержащий речь нескольких спикеров на разных языках. На выходе система должна сформировать структурированную текстовую расшифровку, в которой каждому высказыванию сопоставлены идентификатор спикера, временные метки начала и конца, определённый язык и распознанный текст.

Данная задача естественным образом декомпозируется на четыре последовательных подзадачи:

\begin{enumerate}
    \item \textit{Диаризация спикеров} (speaker diarization) --- определение временных границ речевых фрагментов и их атрибуция конкретным спикерам, то есть ответ на вопрос <<кто говорил и когда>>.
    \item \textit{Сегментация аудио} --- разделение исходного аудиофайла на отдельные аудиофрагменты в соответствии с полученной разметкой диаризации.
    \item \textit{Определение языка} (Language Identification, LID) --- классификация каждого аудиофрагмента по языку произнесённой речи.
    \item \textit{Транскрибация} (Automatic Speech Recognition, ASR) --- преобразование каждого аудиофрагмента в текст с использованием модели, соответствующей определённому языку.
\end{enumerate}

Последовательный характер указанных подзадач определяет выбор pipeline-архитектуры в качестве основного архитектурного паттерна системы~\cite{hohpe2003enterprise}.

К разрабатываемой системе предъявляется ряд ключевых требований, обусловленных предметной областью:

\begin{itemize}
    \item \textit{Локальность обработки}: все данные должны обрабатываться на оборудовании пользователя без передачи в облачные сервисы, что исключает риски утечки конфиденциальной информации.
    \item \textit{Работа на стандартном оборудовании}: система должна функционировать на ПК с 16~ГБ оперативной памяти без необходимости наличия GPU. Допустимое время обработки одного часа аудио составляет до 24 часов.
    \item \textit{Мультиязычность}: поддержка нескольких языков в рамках одной записи, включая языки с ограниченными ресурсами (low-resource languages).
    \item \textit{Расширяемость}: возможность подключения и замены моделей распознавания без модификации основного кода.
\end{itemize}

Для количественной оценки качества работы системы используются стандартные метрики. Для оценки качества распознавания речи применяется метрика WER (Word Error Rate) --- отношение числа ошибочно вставленных, удалённых и заменённых слов к общему числу слов в эталонной расшифровке~\cite{morris2004wer}. Для оценки качества диаризации используется метрика DER (Diarization Error Rate), учитывающая три типа ошибок: пропущенная речь (missed speech), ложное обнаружение речи (false alarm) и ошибка атрибуции спикера (speaker confusion)~\cite{park2022review_der}. Для оценки точности определения языка применяется показатель error rate на уровне сегментов.


\subsection{Обзор методов и моделей}

\subsubsection{Модели автоматического распознавания речи}

Область автоматического распознавания речи переживает период интенсивного развития, связанный с переходом от классических гибридных систем на основе скрытых марковских моделей к сквозным нейросетевым архитектурам. В настоящем обзоре рассматриваются модели, наиболее релевантные для задачи мультиязычной оффлайн-транскрибации.

\paragraph{Whisper.}
Модель Whisper, представленная компанией OpenAI~\cite{radford2023whisper}, является сквозной моделью типа encoder-decoder на основе архитектуры Transformer~\cite{vaswani2017attention}. Whisper обучена на более чем 680\,000 часов слабо размеченных мультиязычных аудиоданных и поддерживает распознавание речи на 99 языках, а также перевод речи и определение языка. На бенчмарке LibriSpeech~\cite{panayotov2015librispeech} модель Whisper large-v3 демонстрирует WER порядка 2,7\% на чистой речи и около 5,2\% в сложных акустических условиях. На мультиязычном бенчмарке Common Voice~15 средний WER составляет около 10\% по всем поддерживаемым языкам. Модель распространяется под лицензией MIT.

Ключевой особенностью Whisper является мультизадачная архитектура: единая модель выполняет распознавание речи, перевод, определение языка и обнаружение голосовой активности, что позволяет заменить несколько компонентов традиционного речевого pipeline~\cite{radford2023whisper}.

Для повышения скорости инференса на CPU разработаны оптимизированные реализации: faster-whisper~\cite{faster_whisper}, использующая движок CTranslate2 для ускорения до~4 раз при сохранении точности, и whisper.cpp~\cite{whisper_cpp} --- порт модели на C/C++ с поддержкой квантизации.

\paragraph{Conformer.}
Архитектура Conformer~\cite{gulati2020conformer}, предложенная исследователями из Google, объединяет механизм самовнимания (self-attention) из Transformer с одномерными свёрточными слоями для совместного моделирования глобальных и локальных зависимостей аудиосигнала. На бенчмарке LibriSpeech Conformer достигает WER 1,9\%~/~3,9\% (test-clean~/~test-other) с внешней языковой моделью. Архитектура Conformer лежит в основе многих современных ASR-систем, включая NVIDIA NeMo и Google USM.

\paragraph{wav2vec 2.0.}
Модель wav2vec~2.0~\cite{baevski2020wav2vec} реализует подход самообучения (self-supervised learning): энкодер предварительно обучается на неразмеченных аудиоданных посредством контрастной задачи, а затем дообучается на небольшом объёме размеченных данных. Данный подход особенно эффективен для языков с ограниченными ресурсами, где объём размеченных данных невелик.

\paragraph{Vosk.}
Vosk~\cite{vosk} --- офлайн-библиотека распознавания речи, поддерживающая около 20 языков. Vosk основан на фреймворке Kaldi и предоставляет компактные модели, пригодные для работы на CPU. Основным преимуществом является малый размер моделей (50--300 МБ) и низкие требования к ресурсам, однако точность распознавания уступает моделям класса Whisper, особенно на мультиязычных данных.

В таблице~\ref{tab:asr_comparison} приведено сравнение рассмотренных ASR-моделей по ключевым характеристикам.

\begin{table}[h]
\centering
\caption{Сравнение моделей автоматического распознавания речи}
\label{tab:asr_comparison}
\small
\begin{tabular}{|l|c|c|c|c|c|}
\hline
\textbf{Модель} & \textbf{Языки} & \textbf{WER (en)} & \textbf{Размер} & \textbf{CPU} & \textbf{Лицензия} \\
\hline
Whisper large-v3 & 99 & 2,7\% & 1,55 ГБ & да* & MIT \\
\hline
Conformer (L) & зав. от обуч. & 1,9\% & $\sim$120М & зав. от реал. & Apache 2.0 \\
\hline
wav2vec 2.0 & зав. от обуч. & 1,8\% & $\sim$300М & да & MIT \\
\hline
Vosk & $\sim$20 & 5--10\% & 50--300 МБ & да & Apache 2.0 \\
\hline
\end{tabular}

\smallskip
\footnotesize{* --- через оптимизированные реализации (faster-whisper, whisper.cpp)}
\end{table}


\subsubsection{Модели диаризации спикеров}

Диаризация спикеров (speaker diarization) --- задача автоматического определения временных границ речевых фрагментов и их атрибуции конкретным спикерам~\cite{park2022review_der}. Современные подходы к диаризации можно разделить на два основных направления: модульные pipeline-системы и сквозные (end-to-end) методы.

\paragraph{pyannote.audio.}
Библиотека pyannote.audio~\cite{bredin2023pyannote, bredin2020pyannote} представляет собой открытый инструментарий на Python для задач диаризации спикеров. Версия~3.1 pipeline диаризации включает три основных этапа:

\begin{enumerate}
    \item \textit{Сегментация спикеров}: нейросетевая модель сегментации применяется к скользящему окну аудио и определяет активность каждого спикера на локальном уровне. Модель обучена с использованием Powerset multi-class cross-entropy loss~\cite{plaquet2023powerset}.
    \item \textit{Извлечение эмбеддингов}: для каждого локально определённого спикера извлекается нейросетевой эмбеддинг.
    \item \textit{Агломеративная кластеризация}: локальные эмбеддинги объединяются в глобальные кластеры, каждый из которых соответствует одному спикеру.
\end{enumerate}

Pipeline pyannote~3.1 демонстрирует DER порядка 18,8\% на датасете AMI~(IHM), 12,2\% на AISHELL-4 и 21,4\% на DIHARD~3 в наиболее строгих условиях оценки (без forgiveness collar, с учётом перекрывающейся речи). Pipeline обучен на комбинации датасетов AISHELL, AliMeeting, AMI, AVA-AVD, DIHARD, Ego4D, MSDWild, REPERE и VoxConverse. Обработка одного часа аудио на GPU (NVIDIA V100) занимает около 1,5~минуты; на CPU время обработки существенно возрастает, однако библиотека полностью поддерживает CPU-инференс.

\paragraph{NVIDIA NeMo.}
Фреймворк NVIDIA NeMo предоставляет альтернативный pipeline диаризации, основанный на модели TitaNet~\cite{snyder2018xvectors} для извлечения эмбеддингов и кластеризации. NeMo ориентирован преимущественно на GPU-инференс, что ограничивает его применимость в контексте требований настоящей работы.

Обоснованием выбора pyannote.audio в качестве основного инструмента диаризации служат: открытый исходный код, поддержка CPU-инференса, высокое качество на стандартных бенчмарках, модульная архитектура и активное сообщество разработчиков.


\subsubsection{Модели определения языка}

Задача определения языка по аудиоданным (Spoken Language Identification, LID) заключается в классификации аудиофрагмента по языку произнесённой речи.

\paragraph{SpeechBrain / ECAPA-TDNN.}
Модель speechbrain/lang-id-voxlingua107-ecapa~\cite{ravanelli2021speechbrain} представляет собой классификатор языка, обученный на наборе данных VoxLingua107~\cite{valk2021voxlingua107} с использованием архитектуры ECAPA-TDNN~\cite{desplanques2020ecapa}. Архитектура ECAPA-TDNN расширяет классическую TDNN (Time Delay Neural Network) механизмами канального внимания (Squeeze-and-Excitation), многоуровневой агрегацией признаков и улучшенным статистическим пулингом. Набор данных VoxLingua107 содержит 6628 часов речи на 107 языках, собранных из видеозаписей YouTube. Модель демонстрирует уровень ошибок 6,7\% на валидационном множестве VoxLingua107.

\paragraph{Whisper LID.}
Модель Whisper~\cite{radford2023whisper} обладает встроенной функцией определения языка: декодер модели включает специальные токены языковых идентификаторов, позволяющие классифицировать язык на этапе инференса. Точность определения языка Whisper коррелирует с объёмом обучающих данных для конкретного языка и составляет 90--99\% для высокоресурсных языков.

В контексте настоящей работы преимуществом модели SpeechBrain является её специализация на задаче LID, более широкий охват языков (107 языков) и возможность использования независимо от ASR-модели, что соответствует принципу модульности разрабатываемой системы.


\subsection{Анализ существующих решений}

Существующие решения для транскрибации многоязычной речи можно разделить на две категории: облачные сервисы и локальные инструменты.

\paragraph{Облачные сервисы.}
К наиболее распространённым облачным решениям относятся:

\begin{itemize}
    \item \textit{Яндекс SpeechKit} --- облачный сервис распознавания и синтеза речи, поддерживающий русский, английский, немецкий, французский и другие языки. Предоставляет API для потоковой и пакетной обработки, а также функцию диаризации. Требует передачи данных в облако Яндекса.
    \item \textit{Сбер SaluteSpeech} --- платформа распознавания речи от Сбера, ориентированная на русский язык. Обеспечивает высокое качество распознавания русской речи, однако мультиязычные возможности ограничены.
    \item \textit{Google Cloud Speech-to-Text} --- облачный сервис, поддерживающий более 125 языков, с функциями диаризации и автоматического определения языка. Является одним из наиболее функциональных решений, однако требует передачи данных в облако Google.
    \item \textit{Stenogram.ai} --- специализированный сервис для расшифровки аудиозаписей встреч с поддержкой диаризации и мультиязычности.
\end{itemize}

Ключевым недостатком всех облачных решений является необходимость передачи аудиоданных на удалённые серверы, что создаёт неприемлемые риски для обработки конфиденциальных записей.

\paragraph{Локальные инструменты.}
Среди локальных решений можно выделить:

\begin{itemize}
    \item \textit{whisper.cpp}~\cite{whisper_cpp} --- реализация модели Whisper на C/C++ с поддержкой CPU-инференса и квантизации. Предоставляет только функцию транскрибации, без диаризации и мультиязычной маршрутизации моделей. Интерфейс ограничен командной строкой.
    \item \textit{Vosk Server}~\cite{vosk} --- серверная реализация офлайн-распознавания на базе Kaldi. Поддерживает около 20 языков, однако не предоставляет диаризацию и определение языка. Модели предустановлены и не допускают замену пользователем.
    \item \textit{faster-whisper}~\cite{faster_whisper} --- оптимизированная реализация Whisper на основе CTranslate2 с ускорением до 4 раз. Как и whisper.cpp, является исключительно ASR-инструментом без интеграции с диаризацией и LID.
\end{itemize}

В таблице~\ref{tab:solutions_comparison} приведено сравнение существующих решений по ключевым критериям, релевантным для задачи настоящей работы.

\begin{table}[h]
\centering
\caption{Сравнение существующих решений для транскрибации речи}
\label{tab:solutions_comparison}
\small
\begin{tabular}{|l|c|c|c|c|c|c|}
\hline
\textbf{Решение} & \textbf{Локальн.} & \textbf{Мультияз.} & \textbf{Диариз.} & \textbf{Расшир.} & \textbf{GUI} & \textbf{CPU} \\
\hline
Яндекс SpeechKit & нет & частично & да & нет & нет & -- \\
\hline
Сбер SaluteSpeech & нет & частично & да & нет & нет & -- \\
\hline
Google Cloud STT & нет & да & да & нет & нет & -- \\
\hline
Stenogram.ai & нет & да & да & нет & да & -- \\
\hline
whisper.cpp & да & да & нет & нет & нет & да \\
\hline
Vosk Server & да & частично & нет & нет & нет & да \\
\hline
faster-whisper & да & да & нет & нет & нет & да \\
\hline
\textbf{Разраб. система} & \textbf{да} & \textbf{да} & \textbf{да} & \textbf{да} & \textbf{да} & \textbf{да} \\
\hline
\end{tabular}
\end{table}

Анализ показывает, что ни одно из существующих решений не объединяет все требуемые характеристики: локальность обработки, полную мультиязычность с маршрутизацией ASR-моделей по языку, диаризацию, расширяемость через plugin-систему, графический интерфейс и работу на CPU. Облачные сервисы решают задачу комплексно, но исключают локальную обработку. Локальные инструменты предоставляют лишь отдельные компоненты (как правило, только ASR), требуя от пользователя самостоятельной интеграции.


\subsection{Выводы по главе и требования к системе}

На основании проведённого анализа предметной области, существующих моделей и программных решений можно сформулировать следующие выводы:

\begin{enumerate}
    \item Задача транскрибации многоязычной речи с разделением по спикерам декомпозируется на четыре подзадачи (диаризация, сегментация, LID, ASR), каждая из которых решается специализированными моделями машинного обучения.
    \item Современные модели ASR (Whisper~\cite{radford2023whisper}), диаризации (pyannote.audio~\cite{bredin2023pyannote}) и определения языка (ECAPA-TDNN на VoxLingua107~\cite{valk2021voxlingua107, desplanques2020ecapa}) достигли уровня качества, достаточного для практического применения.
    \item Существующие решения не объединяют указанные компоненты в единую локальную систему с поддержкой расширяемости и графическим интерфейсом.
\end{enumerate}

На основании анализа сформулированы следующие \textbf{функциональные требования} к разрабатываемой системе:

\begin{enumerate}
    \item загрузка аудиофайлов распространённых форматов (WAV, MP3, FLAC, OGG);
    \item автоматическая диаризация спикеров с использованием pyannote.audio~\cite{bredin2023pyannote};
    \item сегментация аудиозаписи на отдельные фрагменты в соответствии с разметкой диаризации;
    \item определение языка каждого аудиофрагмента с использованием модели LID~\cite{valk2021voxlingua107};
    \item транскрибация каждого фрагмента с автоматическим выбором ASR-модели в зависимости от определённого языка;
    \item формирование итоговой расшифровки, структурированной по спикерам с метками языка и временными метками;
    \item экспорт результатов в форматы JSON и TXT;
    \item графический интерфейс с возможностью загрузки файла, запуска обработки, отображения прогресса и просмотра результатов.
\end{enumerate}

\textbf{Нефункциональные требования} к разрабатываемой системе:

\begin{enumerate}
    \item полностью локальная обработка без передачи данных в облачные сервисы;
    \item работа на стандартном оборудовании: CPU, 16~ГБ ОЗУ, без GPU;
    \item время обработки одного часа аудио не должно превышать 24~часов на CPU;
    \item кроссплатформенность: поддержка операционных систем Windows, Linux, macOS;
    \item воспроизводимость: сохранение конфигурации каждого запуска, промежуточных результатов и возможность возобновления прерванной обработки;
    \item модульность: возможность подключения и замены моделей ASR, LID и диаризации через систему подключаемых модулей без модификации основного кода~\cite{gamma1995patterns}.
\end{enumerate}